\documentclass[runningheads,a4paper]{article}
\usepackage[utf8]{inputenc}
\setcounter{tocdepth}{3}
\usepackage[english]{babel} 
\usepackage{graphicx}
\usepackage{grffile}
\usepackage{float}
\usepackage{multicol}
\usepackage{url}
\usepackage{titling}
\usepackage[hidelinks]{hyperref}
\setcounter{secnumdepth}{5}
%Margins
\usepackage[
margin=2cm,
includefoot
]{geometry}
\graphicspath{{img/}}
%Headers and Footers
\usepackage{fancyhdr}
\pagestyle{fancy}
\fancyhead{}
\fancyfoot{}
\fancyfoot[R]{\thepage}
\renewcommand{\headrulewidth}{0pt}
\renewcommand{\footrulewidth}{0pt}
\setlength\parindent{24pt}
\begin{document}
%Title Page
\begin{titlepage}
\begin{center}
\includegraphics[width=10cm]{UP.jpg}  \\
[1cm]
\line(1,0){300} \\
[0.3cm]
\textsc{\Large
NavUP \\
Testing Phase \\
\hfill \break 1 May 2017
%University of Pretoria
}\\
[0.1cm]
\line(1,0){300} \\
[0.7cm]
\textsc{\Large
Team GladIOS Navigation
} \\
\end{center}
\begin{center}
\begin{multicols}{2}
\textsc{\large\\
Jacobus Marais\\ 
u15188397\\ 
}
\textsc{\large\\
Mia Gerber\\
u15016502\\ 
}
\textsc{\large\\
Victor Twigge\\
u10376802\\ 
}
\columnbreak
\textsc{\large\\
Nicaedin Suklal\\
u15207812\\
}
\textsc{\large\\
Darren Adams\\
u14256232\\
}
\textsc{\large\\
Nkosinathi Mothoa\\
u12077420\\
}
\end{multicols}
\textsc{	\\ \href{https://github.com/KobusMarais/cos301_Longsword_Navigation}{GitHub}
\url{https://github.com/KobusMarais/cos301_Longsword_Navigation}}
\end{center}
\end{titlepage}

%\maketitle



\begingroup



\tableofcontents

\addcontentsline{toc}{section}{Table Of Contents}

\endgroup

\newpage

\section{Introduction}

	\section{Non-Functional Requirements}
	
	This section discusses and rates 10 non-functional requirements of the NavUP system.
	Below is the analysis of the Navigation module Interface and all its component classes:
	
	\begin{itemize}
		\item \textbf{Documentation:}
			The module as a whole is well documented. The relavent documentation tools were used effectively. [\textbf{Mark: 10/10}]
			
		\item \textbf{Efficiency:} The module is as efficient as it can be with the current tools and resources used. [\textbf{Mark: 8/10}]
		
		\item \textbf{Maintainability:} The project has good module cohesion, as a result it is easily maintainable. [\textbf{Mark: 8/10}]
		
		\item \textbf{Quality:} The module follows a consistent coding standard throughout all its component classes. The code itself is readable with little effort. [\textbf{Mark: 7/10}]
		
		\item \textbf{Reuseablity:} Due to the explicit requirements of NavUP, not all of the modules are easily reuseable. However slight tweaks to some modules is enough to be able to reuse them effectively. [\textbf{Mark: 5/10}]
		
		\item \textbf{Testability:} description
		\item \textbf{Complaince:} description
		\item \textbf{Robustness:} description
		\item \textbf{Cohesion:} description
		\item \textbf{Coupling:} description\\
		
	\end{itemize}

\end{document}
